\section{Macro-level Consequences}

In this section, I will describe the macro-level consequences of the micro-level assumptions regarding households and landlords discussed earlier as I have implemented them in the agent-based model. A thorough analysis of this computational model can be found in the supplementary material. Here, I will only present results for selected combinations of model parameters, often only the empirically most plausible case: households value both neighbor's status and housing quality but housing quality is more important, income and status are correlated but not perfectly so, and housing units retain most of their quality over short periods even without investments on behalf of landlords\footnote{a = 0.25, r = 0.5, b = 0.95. For a description of these parameters and the other (constant) model parameters, see table \ref{tab:parameters}.}. The results are based on simulation runs conducted after the burn-in period, when the system has reached equilibrium.

\begin{figure}
\caption{Segregation Levels for selected Model Parameters}
\includegraphics[width=\textwidth]{/home/malte/Projects/SocEconHousing/images/abm_segregation}

\tiny{Segregation indices calculated for 5x5 neighborhoods and selected model parameters. Each dot represents the average segregation index for these parameter combinations over 15 independent runs and 100 time steps. The figure is based on simulated data from the main computational experiment (see Appendix A) with the following selection: $b = 0.95$ and $r \in \{0, 0.5, 1\}$.} 
\label{fig:seg}
\end{figure}


Crucially, the model predicts a sizeable amount of residential segregation and spatial inequality in the empirically most likely case, as shown in the highlighted area of Figure \ref{fig:seg}. Rents and housing quality are unevenly distributed across neighbourhoods in this case, while income segregation is less pronounced. Residential segregation by status is much lower but still significantly more uneven than expected by random chance. Based on the full analysis of the model, I infer that cities self-organize into distinct and stable neighborhoods via the following mechanism: When a neighborhood increases in status, it also becomes more desirable. As desirability increases, rents increase, which triggers landlords to invest. Because the attractiveness of the neighbourhood rises with both status and housing quality, rents increase further, and the poorest residents of the neighbourhood are displaced (or poor residents can no longer afford to move to the neighbourhood). Because low incomes correlate with low status, the residents with the lowest status are likely to move out, which increases the average status of the neighbourhood. This self-reinforcing dynamic continues until the model reaches an equilibrium, where a marginal increase in neighborhood status and/or housing quality is unable to attract wealthier residents from other neighborhoods. This dynamic aligns well with the reviewed descriptions of gentrification. Conversely, if a neighbourhood decreases in status, rents tend to decrease, and landlords are less likely to invest in it. As neighborhood status and housing quality decrease, affluent households move to different neighborhoods where they can afford a higher residential satisfaction. The decreasing rent attracts less wealthy residents, for whom it is still a relatively good neighborhood. 

\begin{figure}
\caption{Residential Mobility by Income Deciles}
\includegraphics[width=\textwidth]{/home/malte/Projects/SocEconHousing/images/abm_mobility}

\tiny{Frequencies of the residential status from t to t+1. If a household occupies the same housing unit in t+1 as in t, the household stayed. If the household changed the housing unit, they moved. If at t, the rent of the unit was greater than the household's disposable income, they are required to move (except if they already live in the cheapest available housing unit). If the household moved and the rent was unaffordable at t, the move is considered a forced move. The first decile has the lowest incomes, the tenth decile the highest. The figure is based on simulated data from the main computational experiment (see Appendix A) with the following selection: $a = 0.25$, $r = 0.5$ and $b = 0.95$.}
\label{fig:mobility}
\end{figure}

These up- and downward spirals of neighborhoods continue until the overall city reaches an equilibrium. Once the city reaches such a stratification, where the rich reside in the best housing units in high-status neighborhoods, the richest have no incentive to move to different neighborhoods and housing units, as they cannot improve their residential satisfaction. In turn, as the most desirable places are occupied by the richest and have very high rents, the slightly less affluent households cannot improve their housing situation either and stay in the second-best houses and neighborhoods, etc. Figure \ref{fig:mobility} shows that, in equilibrium, households with high incomes are very likely to stay in the same housing unit from one point in time to the next. Households with low incomes, on the other hand, tend to move more frequently. They often face unaffordable rents, even though they already live in low-quality housing units and neighborhoods (see also figure \ref{fig:quality}). When they cannot afford the rent, due to the model's design, they need to move unless they are already at the cheapest location currently available. Therefore, they move often even though their housing situation is unlikely to improve after the move. Middle-income households also exhibit lower but still considerable levels of residential mobility. However, more often than not, this move is voluntary because an affordable vacancy became available, providing them with greater residential satisfaction. Because my model incorporates population dynamics, households are randomly removed and added to the simulation, allowing vacancies to open up randomly. Because most households have a middle income, and therefore most neighborhoods are in the middle of their respective distributions, most vacancies are both available and attractive to middle-income households. The random creation of vacancies and the utility maximization assumption, which makes households move even for a marginal increase in utility, leads to higher levels of residential mobility for middle-income households in my model than we would expect in the real world. Without population turnover, middle-income households would also move rarely (see robustness check B in the supplement). Although households typically do not move to improve their residential satisfaction marginally, but due to life-course events, the model can reproduce the stylized fact that poor households have higher residential mobility due to affordability constraints, while middle- and high-income households move less often. If they do, they typically are not forced to do so and can improve their residential outcomes. Furthermore, the high residential mobility of low-income households, even in equilibrium, fits well with the studies that find little increased mobility of low-income households in gentrifying neighborhoods \citep{freemanDisplacementSuccessionResidential2005, dingGentrificationResidentialMobility2016, leeGeographyGentrificationResidential2023}. Despite individual residential mobility, neighborhoods are relatively stable in the empirically most plausible condition. However, in line with the stylized facts again, there is an inverse-U-shaped relationship between the variation in income ranks and the average income rank: neighborhoods at the top and bottom of the income distribution tend to be more stable than those in the middle (see supplements A and B).

\begin{figure}
\caption{Household Income, Rent, and Housing Quality}
\includegraphics[width=\textwidth]{/home/malte/Projects/SocEconHousing/images/abm_inequality}

\tiny{LOESS smoothers for the relationship between housing quality, rent, and the rent-to-income ratio with household income in the empirically most plausible parameter combination. The figure is based on simulated data from the main computational experiment (see Appendix A) with the following selection: $a = 0.25$, $r = 0.5$ and $b = 0.95$.} 
\label{fig:quality}
\end{figure}

As discussed in the mechanism leading to segregation, high-income households can outcompete other households for the most desirable housing units. Figure \ref{fig:quality} confirms this relationship on the individual level: there is a positive relationship between household income and both the rent they pay for the housing unit and the quality of the housing unit. However, the slope of this relationship decreases as income increases. While the absolute rent and housing quality (and neighborhood status, see robustness check C) increase with income, the relative increase slows down. The result is the observation that wealthier households pay more in absolute terms but a lower proportion of their income for housing. In contrast, it is common for low-income households to live in housing units where rent exceeds their disposable income, even though they live in cheap, low-quality neighborhoods and housing units. But why are the rents of the worst housing units still higher than the disposable incomes of the poorest? As rent is set globally based on utility, there is no exploitation on behalf of landlords in this model: they do not set prices. The reason is easiest to understand when considering neighborhood status: because neighborhood status is, by definition, an average of the individual status of residents in a neighborhood, the distribution of neighborhood status has lower variance than the individual status distribution. Because landlords invest in housing quality based on the average rent in the neighborhood, there is a reduction in variance in housing quality compared to rents as well. So, for both components of the utility function, there is a reduction in variance because they depend on averages. The distribution of housing units by utility is less unequal than the income distribution of households. For wealthy households, this limits the supply of desirable units, which is a reason why the rich pay a lower share of their income for housing\footnote{This result could also be due to the specification of the utility provided by a unit (Appendix A, equation \ref{eq:util}). It is a common assumption in economics that marginal utility is decreasing. When a household already resides in a high-quality housing unit, the added utility provided by a fixed increase in housing quality is lower than the added utility of the same change in housing quality when the household lives in a housing unit of lower quality (and the same holds for neighborhood status). As a consequence, households are less willing to pay more rent for improved housing quality and a better neighborhood status. Moreover, the choice of using the 75th percentile in the rent formula could contribute to the slowing down of the relationship between household income and utility.}. However, for the poorest, this means that no housing units are supplied that are of such low quality that they could afford them. 

\begin{figure}
\caption{Regression of Housing Quality on Status and Income}
\includegraphics[width=\textwidth]{/home/malte/Projects/SocEconHousing/images/abm_statusreg}

\tiny{Coefficient of household status with 99.9\% confidence intervals from separate OLS regressions of housing quality on household income and status. Each regression is calculated using data from all runs and time points with the same model parameters. Because the agents in the model are interdependent and observed over time, the IID assumption of the OLS regression is violated, resulting in downward-biased standard errors. The figure is based on simulated data from the main computational experiment (see Appendix A) with the following selection: $r = 0.5$ and $b \in \{0, 0.95, 1\}$. I also omitted the case of $a = 1$ because this analysis does not make much sense if households do not have a preference for housing quality.}
\label{fig:reg}
\end{figure}

Because status and income are correlated, it is no surprise that low-status households tend to reside in poorer neighborhoods and substandard housing units. However, existing empirical research has shown that minorities often have worse residential outcomes even after controlling for their income. As discussed above, landlords cannot set prices in this model and, therefore, this model cannot reproduce price discrimination. Nevertheless, this model does contain a plausible mechanism why lower status can be associated lower with housing quality independent of income. Figure \ref{fig:reg} shows the OLS coefficients for status from multiple regressions of housing quality on household status and income. We observe a significantly positive relationship between status and housing quality after controlling for income, which increases as households value neighborhood status more in their location choice. High-status households tend to live in higher-quality housing units, and low-status households live in lower-quality housing units, independent of income. The reason for this observation within this model is that when a low-status household moves to a neighborhood, the attractiveness of that neighborhood and, consequently, its rents decrease. When rents decrease, landlords tend to refrain from investing, resulting in a deterioration of housing quality. The same mechanism holds in reverse: a high-status household increases neighborhood attractiveness and rents, which leads landlords to invest in a higher housing quality. This result aligns well with empirical studies on racial discrimination in the US, where a frequent concern among existing residents is that housing values will decrease when minority members move to the neighborhood \citep{harrisPropertyValuesDrop1999, krysanDoesRaceMatter2009}. However, this relationship is moderated by the level of inertia in housing quality. With high inertia, housing quality reacts very slowly to changes in local population and changes are small. In the highlighted empirically most likely case, the effect is even negative with high inertia. Due to deviation from the overall pattern and the likely underestimation of standard errors (see figure description), I consider these null effects though. Whether this theoretically plausible mechanism therefore persist in the real world with a high inertia in housing quality, is an empirical question.

% TODO add vacancies?